% =====================================================================
%  BOZZA - Capitolo 5 (italiano)
%  Conclusioni: riepilogo del lavoro di tirocinio, risultati principali,
%  limiti dello studio, sviluppi futuri e considerazioni finali.
% =====================================================================
\chapter{Conclusioni}
\label{ch:conclusioni}

Questo capitolo chiude la tesi ripercorrendo il lavoro svolto durante il
tirocinio e i suoi esiti. Si riepilogano dapprima le attivit\`a condotte
rispetto agli obiettivi assegnati (Sezione~\ref{sec:concl-riepilogo}), quindi
i risultati principali della valutazione sperimentale
(Sezione~\ref{sec:concl-risultati}). Seguono i limiti dello studio
(Sezione~\ref{sec:concl-limiti}), le direzioni di sviluppo che il lavoro
lascia aperte (Sezione~\ref{sec:concl-futuro}) e alcune considerazioni
conclusive (Sezione~\ref{sec:concl-finale}).

% ---------------------------------------------------------------------
\section{Riepilogo del lavoro svolto}
\label{sec:concl-riepilogo}
Il tirocinio si \`e inserito nel progetto ProSafe~\cite{prosafe}, il sistema
di boe intelligenti per la sicurezza di prossimit\`a nelle acque costiere
sviluppato dal gruppo di ricerca ospitante. Il protocollo di partenza limita
la conoscenza di ogni nodo ai vicini entro la portata radio; l'obiettivo del
progetto era estendere questa conoscenza oltre il primo hop, contenendone il
costo sul canale, e quantificarne il compromesso con una valutazione
sperimentale sistematica (Sezione~\ref{sec:intro-obiettivi}).

Tutti e tre gli obiettivi assegnati sono stati raggiunti. Sono state
progettate due modalit\`a di disseminazione multi-hop dagli approcci opposti:
il \emph{forwarding aggregato}, che allega a ogni beacon la lista dei nodi
conosciuti senza generare frame aggiuntivi, e il \emph{forwarding singolo},
che rilancia i beacon ricevuti entro un limite di hop con meccanismi di
contenimento del broadcast storm (Capitolo~\ref{ch:contributo}). Entrambe
sono state implementate nel simulatore a eventi discreti del progetto,
integrandole nella pipeline di trasmissione CSMA/CA esistente senza alterare
il comportamento delle componenti preesistenti, e affiancandovi le metriche
necessarie a misurarne gli effetti (Capitolo~\ref{ch:simulatore}). Le
modalit\`a sono state infine confrontate in una campagna che combina due
condizioni di canale, tre intervalli di beaconing e tre protocolli di
scheduling su cinque livelli di densit\`a, con 20 repliche indipendenti per
configurazione (Capitolo~\ref{ch:risultati}).

% ---------------------------------------------------------------------
\section{Risultati principali}
\label{sec:concl-risultati}
La valutazione sperimentale ha prodotto quattro indicazioni principali.

\paragraph{L'inoltro multi-hop \`e una necessit\`a, non un'ottimizzazione.}
Nello scenario adottato --- area ampia rispetto alla portata radio e
popolazione in continuo ricambio --- il beaconing single-hop si ferma al
63--77\% di rete scoperta nel canale ideale e al 50--68\% in quello con
perdite, senza mai convergere. Entrambe le modalit\`a multi-hop completano
invece di fatto la scoperta dai 60 nodi in su, in entrambe le condizioni di
canale.

\paragraph{Il forwarding aggregato domina il forwarding singolo.} A parit\`a
di copertura, l'aggregato paga una frazione del costo in contesa del singolo
(a 100 boe nel canale ideale, collision rate 0.163 contro 0.553 con SBP e
0.035 contro 0.421 con ACAB) e diffonde la conoscenza pi\`u rapidamente
(latenza di reazione di circa \SI{16}{ms} contro 89--115~\si{ms}): affidare
la disseminazione all'allungamento dei frame anzich\'e a trasmissioni
aggiuntive si \`e rivelata la scelta pi\`u efficiente lungo tutti gli assi
misurati. Il forwarding singolo conserva la semplicit\`a del meccanismo e un
delivery ratio migliore nei soli scenari sparsi con canale con perdite.

\paragraph{Il beaconing adattivo rende l'inoltro sostenibile.} ADAB e ACAB
assorbono quasi interamente la contesa aggiuntiva introdotta dal forwarding
aggregato --- mantenendo il B-PDR sopra 0.94 nel canale ideale fino a 100 boe
--- e rendono le prestazioni insensibili all'intervallo di beaconing
configurato, l\`a dove SBP degrada rapidamente con densit\`a e frequenza di
trasmissione. Segnali puramente locali, come il conteggio dei vicini,
bastano dunque a regolare il traffico anche quando il carico \`e amplificato
dall'inoltro, senza alcun coordinamento centralizzato.

\paragraph{La combinazione migliore dipende dal canale.} Nelle condizioni
realistiche di canale con perdite, ADAB--forwarding aggregato \`e emerso come
il compromesso complessivamente migliore: scoperta pressoch\'e completa a
ogni densit\`a e affidabilit\`a al livello del semplice single-hop.
ACAB--forwarding aggregato massimizza invece l'affidabilit\`a dove il canale
\`e buono, ma il suo backoff pi\`u conservativo rallenta la scoperta negli
scenari sparsi con perdite, dove la ridondanza delle trasmissioni \`e
proprio ci\`o da cui l'inoltro trae robustezza.

% ---------------------------------------------------------------------
\section{Limiti dello studio}
\label{sec:concl-limiti}
I risultati vanno letti entro i confini della campagna condotta. L'intervallo
di densit\`a esplorato (20--100 boe, da circa 2 a circa 9 vicini medi) \`e
rappresentativo di scenari costieri realistici, ma le tendenze osservate
indicano che il forwarding singolo con SBP \`e gi\`a prossimo alla
saturazione al limite superiore, e la crescita monotona del collision rate
suggerisce che nessuna combinazione statica reggerebbe densit\`a molto
superiori; per il forwarding aggregato il fattore limitante di lungo periodo
\`e la dimensione del beacon, che cresce con il numero di nodi annunciati e
allunga i tempi di occupazione del canale. Il modello di canale, pur
calibrato su misure sul campo, resta una semplificazione a due soglie della
propagazione reale, e il modello di mobilit\`a Random Waypoint approssima ma
non riproduce il comportamento di nuotatori e piccoli mezzi. Infine, la
valutazione \`e interamente simulativa: i risultati quantificano i rapporti
tra le alternative progettuali nelle stesse identiche condizioni, ma i valori
assoluti andranno confermati su hardware reale.

% ---------------------------------------------------------------------
\section{Sviluppi futuri}
\label{sec:concl-futuro}
Il lavoro lascia aperte alcune direzioni, in parte gi\`a predisposte nel
codice durante il tirocinio.

\paragraph{Attivazione dei meccanismi progettati.} I due accorgimenti
descritti nella Sezione~\ref{sec:improvements} sono presenti nel simulatore
ma non sono stati impiegati negli esperimenti: il \emph{gate di forwarding
density-aware} (Eq.~\ref{eq:fwdgate}), che dirada probabilisticamente gli
inoltri nelle zone dense, potrebbe ridurre il divario di collision rate che
penalizza il forwarding singolo; il \emph{limite di hop per il forwarding
aggregato} conterrebbe invece la crescita della dimensione del beacon,
il fattore limitante della modalit\`a alle densit\`a estreme. La loro
valutazione e taratura \`e il seguito pi\`u naturale di questo lavoro.

\paragraph{Regimi di densit\`a superiori.} Estendere la scala oltre le 100
boe permetterebbe di osservare direttamente i punti di saturazione anticipati
dalle tendenze misurate, e di verificare fino a dove il backoff adattivo
riesce a rinviare la saturazione dell'airtime del forwarding aggregato.

\paragraph{Ottimizzazione di ACAB per canali con perdite.} La lentezza di
scoperta di ACAB--aggregato negli scenari sparsi con perdite suggerisce di
rivederne pesi e soglie --- o di renderli adattivi rispetto alla qualit\`a
del canale percepita --- cos\`i da recuperare ridondanza dove serve senza
rinunciare alla prudenza dove il canale \`e affidabile.

\paragraph{Riduzione della latenza di reazione del forwarding singolo.} La
coda di inoltro FIFO tratta tutte le origini allo stesso modo; politiche di
priorit\`a (ad esempio privilegiare le origini pi\`u lontane o meno
annunciate) potrebbero abbattere la latenza di reazione, il principale
svantaggio temporale della modalit\`a.

\paragraph{Validazione sul campo.} Il passo conclusivo \`e il confronto dei
risultati simulativi con esperimenti su hardware reale nell'ambito del
progetto ProSafe, includendo aspetti che il simulatore non modella: consumo
energetico delle boe, interferenza esterna sul canale e comportamento dei
chipset Wi-Fi commerciali in trasmissione broadcast.

% ---------------------------------------------------------------------
\section{Considerazioni conclusive}
\label{sec:concl-finale}
Sul piano del progetto, il tirocinio ha consegnato al gruppo di ricerca un
simulatore esteso con due modalit\`a di disseminazione multi-hop
configurabili, una suite di metriche pi\`u ricca e una campagna sperimentale
riproducibile che quantifica il compromesso tra copertura, affidabilit\`a e
scalabilit\`a: elementi che potranno sostenere le evoluzioni future di
ProSafe. Sul piano formativo, il lavoro ha richiesto di attraversare l'intero
ciclo di un'attivit\`a di ricerca --- lo studio di un sistema esistente, la
progettazione di meccanismi nuovi sotto vincoli stringenti, l'implementazione
rigorosa in un'architettura a eventi discreti e l'analisi statistica dei
risultati --- confermando quanto la simulazione, se validata e usata con
consapevolezza dei suoi limiti, sia uno strumento essenziale per lo studio
dei protocolli di rete l\`a dove la sperimentazione diretta \`e
impraticabile.
